% Part: computability
% Chapter: computability-theory
% Section: s-m-n

\documentclass[../../../include/open-logic-section]{subfiles}

\begin{document}

\olfileid[fr]{cmp}{thy}{smn}
\olsection{Le théorème $s$-$m$-$n$}

\begin{explain}
Le théorème suivant porte le nom de « théorème
$s$-$m$-$n$ », pour une raison qui apparaîtra
bientôt. La difficulté principale est de bien
comprendre son énoncé; une fois celui-ci compris,
le résultat semblera assez naturel.
\end{explain}

\begin{thm}
\ollabel{thm:s-m-n}
  Pour chaque paire d'entiers naturels $n$ et~$m$,
  il existe une fonction récursive primitive~$s^m_n$
  telle que, pour toute suite d'entiers
  $e$, $a_0$, \dots, $a_{m-1}$,
  $y_0$, \dots, $y_{n-1}$, on ait
  \[
  \cfind{s^m_n(e, a_0, \dots, a_{m-1})}[n](y_0, \dots, y_{n-1}) \simeq
  \cfind{e}[m+n](a_0, \dots, a_{m-1}, y_0, \dots, y_{n-1}).
\]
\end{thm}

\begin{explain}
Il est utile de considérer $s^m_n$ comme une
fonction agissant sur des \emph{programmes}.
Elle reçoit un programme~$e$ calculant une
fonction à $(m+n)$~arguments, ainsi que
les entrées fixées $a_0,\dots,a_{m-1}$,
et renvoie un programme de code
$s^m_n(e,a_0,\dots,a_{m-1})$ qui calcule
la fonction des $n$~arguments restants. Si $m=0$, aucune entrée n'est fixée ; si $n=0$, on admet ici les programmes sans entrée, qui produisent une valeur ou divergent.
\iftag{TMs}{Si $e$ décrit une machine de Turing,
$s^m_n(e,a_0,\dots,a_{m-1})$ est le code d'une
machine qui, sur l'entrée $y_0,\dots,y_{n-1}$,
place d'abord $a_0,\dots,a_{m-1}$ en tête
de la suite d'entrées, puis exécute~$e$.
Chaque $s^m_n$ est alors une fonction
récursive primitive qui produit le code
de la machine correspondante.}{}\footnote{Corrections éditoriales : dans l'explication source, deux occurrences de $x$ sont remplacées par l'indice de programme $e$ déjà fixé. La valeur de $s^m_n$ est un code de machine, plutôt que la machine elle-même.}
\end{explain}

\end{document}
